% ======================================================================
% Chapter 7: Radiative Corrections: Some Formal Developments
% (book pages 211-257; PDF pages 234-287)
% ======================================================================
\chapter{Radiative Corrections: Some Formal Developments}

We cheated four times in the last three chapters,%
\footnote{A fifth cheat, postulating rather than deriving the photon
propagator, will be remedied in Chapter~9.} stating (and sometimes
motivating) a result but postponing its proof.  Those results were:
\begin{enumerate}
  \item The formula for decay rates in terms of $S$-matrix elements,
    Eq.~(4.86).
  \item The master formula for $S$-matrix elements in terms of Feynman
    diagrams, Eq.~\eqref{eq:4.103}.
  \item The Ward identity, Eq.~\eqref{eq:5.79}.
  \item The ad hoc subtraction to remove the ultraviolet divergence in the
    vertex-correction diagram, Eq.~(6.55).
\end{enumerate}
It is time now to return to these issues and give them a proper treatment.
In Sections~7.2 through~7.4 we will derive all four of these results.  The
knowledge we gain along the way will help us interpret the three remaining
loop corrections for electron scattering from a heavy target shown in
\eqref{eq:6.1}: the external leg corrections and the vacuum polarization.  We
will evaluate the former in Section~7.1 and the latter in Section~7.5.

This chapter will be more abstract than the two preceding ones.  Its main
theme will be the singularities of Feynman diagrams viewed as analytic
functions of their external momenta.  We will find, however, that this
apparently esoteric subject is rich in physical implications, and that it
illuminates the relation between Feynman diagrams and the general principles
of quantum theory.

\section{Field-Strength Renormalization}

In this section we will investigate the analytic structure of the two-point
correlation function,
\begin{equation}
  \langle\Omega|T\phi(x)\phi(y)|\Omega\rangle
  \qquad\text{or}\qquad
  \langle\Omega|T\psi_\alpha(x)\bar{\psi}_\beta(y)|\Omega\rangle .
\end{equation}
In a free field theory, the two-point function
$\langle 0|T\phi(x)\phi(y)|0\rangle$ has a simple interpretation: It is the
amplitude for a particle to propagate from $y$ to $x$.  To what extent does
this interpretation carry over into an interacting theory?

Our analysis of the two-point function will rely only on general principles
of relativity and quantum mechanics; it will not depend on the nature of the
interactions or on an expansion in perturbation theory.  We will, however,
restrict our consideration to scalar fields.  Similar results can be obtained
for correlation functions of fields with spin; we will display the analogous
result for Dirac fields at the end of the analysis.

To dissect the two-point function $\langle\Omega|T\phi(x)\phi(y)|\Omega\rangle$
we will insert the identity operator, in the form of a sum over a complete set
of states, between $\phi(x)$ and $\phi(y)$.  We choose these states to be
eigenstates of the full interacting Hamiltonian, $H$.  Since the momentum
operator $\mathbf{P}$ commutes with $H$, we can also choose the states to be
eigenstates of $\mathbf{P}$.  But we can also make a stronger use of Lorentz
invariance.  Let $|\lambda_0\rangle$ be an eigenstate of $H$ with momentum
zero: $\mathbf{P}|\lambda_0\rangle = 0$.  Then all the boosts of
$|\lambda_0\rangle$ are also eigenstates of $H$, and these have all possible
3-momenta.  Conversely, any eigenstate of $H$ with definite momentum can be
written as the boost of some zero-momentum eigenstate $|\lambda_0\rangle$.
The eigenvalues of the 4-momentum operator $P^\mu = (H, \mathbf{P})$ organize
themselves into hyperboloids, as shown in Fig.~\ref{fig:7.1}.

\begin{figure}[t]
  \centering
  \includegraphics[width=0.55\textwidth]{fig7.1.png}
  \caption{The eigenvalues of the 4-momentum operator
    $P^\mu = (H, \mathbf{P})$ occupy a set of hyperboloids in energy-momentum
    space.  For a typical theory the states consist of one or more particles of
    mass $m$.  Thus there is a hyperboloid of one-particle states and a
    continuum of hyperboloids of two-particle states, three-particle states,
    and so on.  There may also be one or more bound-state hyperboloids below
    the threshold for creation of two free particles.}
  \label{fig:7.1}
\end{figure}

Recall from Chapter~2 that the completeness relation for the one-particle
states is
\begin{equation}
  (\text{identity})_{\text{1-particle}} = \int
    \frac{d^3p}{(2\pi)^3}\ \frac{1}{2E_{\mathbf{p}}}\ |\mathbf{p}\rangle
    \langle\mathbf{p}| .
  \label{eq:7.1}
\end{equation}
We can write an analogous completeness relation for the entire Hilbert space
with the aid of a bit of notation.  Let $|\lambda_{\mathbf{p}}\rangle$ be the
boost of $|\lambda_0\rangle$ with momentum $\mathbf{p}$, and assume that the
states $|\lambda_{\mathbf{p}}\rangle$, like the one-particle states
$|\mathbf{p}\rangle$, are relativistically normalized.  Let
$E_{\mathbf{p}}(\lambda) = \sqrt{|\mathbf{p}|^2 + m_\lambda^2}$, where
$m_\lambda$ is the ``mass'' of the states $|\lambda_{\mathbf{p}}\rangle$,
that is, the energy of the state $|\lambda_0\rangle$.  Then the desired
completeness relation is
\begin{equation}
  1 = |\Omega\rangle\langle\Omega| + \sum_\lambda \int
    \frac{d^3p}{(2\pi)^3}\ \frac{1}{2E_{\mathbf{p}}(\lambda)}\
    |\lambda_{\mathbf{p}}\rangle\langle\lambda_{\mathbf{p}}|,
  \label{eq:7.2}
\end{equation}
where the sum runs over all zero-momentum states $|\lambda_0\rangle$.

We now insert this expansion between the operators in the two-point function.
Assume for now that $x^0 > y^0$.  Let us drop the uninteresting constant term
$\langle\Omega|\phi(x)|\Omega\rangle\langle\Omega|\phi(y)|\Omega\rangle$.
(This term is usually zero by symmetry; for higher-spin fields, it is zero by
Lorentz invariance.)  The two-point function is then
\begin{equation}
  \langle\Omega|\phi(x)\phi(y)|\Omega\rangle
  = \sum_\lambda \int \frac{d^3p}{(2\pi)^3}\ \frac{1}{2E_{\mathbf{p}}(\lambda)}\
    \langle\Omega|\phi(x)|\lambda_{\mathbf{p}}\rangle\
    \langle\lambda_{\mathbf{p}}|\phi(y)|\Omega\rangle .
  \label{eq:7.3}
\end{equation}
We can manipulate the matrix elements as follows:
\begin{equation}
  \begin{aligned}
    \langle\Omega|\phi(x)|\lambda_{\mathbf{p}}\rangle
    &= \langle\Omega|e^{iP\cdot x}\phi(0)e^{-iP\cdot x}|\lambda_{\mathbf{p}}\rangle\\
    &= \langle\Omega|\phi(0)|\lambda_{\mathbf{p}}\rangle\ e^{-ip\cdot x}
      \Big|_{p^0=E_{\mathbf{p}}}\\
    &= \langle\Omega|\phi(0)|\lambda_0\rangle\ e^{-ip\cdot x}
      \Big|_{p^0=E_{\mathbf{p}}} .
  \end{aligned}
  \label{eq:7.4}
\end{equation}
The last equality is a result of the Lorentz invariance of $|\Omega\rangle$
and $\phi(0)$: Insert factors of $U$, where $U$ is the unitary operator that
implements a Lorentz boost from $\mathbf{p}$ to 0, and use
$U\phi(0)U^{-1} = \phi(0)$.  (For a field with spin, we would need to keep
track of its nontrivial Lorentz transformation.)  Introducing an integration
over $p^0$, our expression for the two-point function (still for $x^0 > y^0$)
becomes
\begin{equation}
  \langle\Omega|\phi(x)\phi(y)|\Omega\rangle
  = \sum_\lambda \int \frac{d^4p}{(2\pi)^4}\ \frac{i}{p^2-m_\lambda^2+i\epsilon}\
    e^{-ip\cdot(x-y)}\ \bigl| \langle\Omega|\phi(0)|\lambda_0\rangle
      \bigr|^2 .
  \label{eq:7.5}
\end{equation}
Note the appearance of the Feynman propagator, $D_F(x-y)$, but with $m$
replaced by $m_\lambda$.

\begin{figure}[t]
  \centering
  \includegraphics[width=0.5\textwidth]{fig7.2.png}
  \caption{The spectral function $\rho(M^2)$ for a typical interacting field
    theory.  The one-particle states contribute a delta function at $m^2$ (the
    square of the particle's mass).  Multiparticle states have a continuous
    spectrum beginning at $(2m)^2$.  There may also be bound states.}
  \label{fig:7.2}
\end{figure}

Analogous expressions hold for the case $y^0 > x^0$.  Both cases can be
summarized in the following general representation of the two-point function
(the K\"all\'en--Lehmann spectral representation):
\begin{equation}
  \langle\Omega|T\phi(x)\phi(y)|\Omega\rangle
  = \int_0^{\infty} \frac{dM^2}{2\pi}\ \rho(M^2)\ D_F(x-y; M^2),
  \label{eq:7.6}
\end{equation}
where $\rho(M^2)$ is a positive spectral density function,
\begin{equation}
  \rho(M^2) = \sum_\lambda (2\pi)\,\delta(M^2 - m_\lambda^2)\,
    \bigl| \langle\Omega|\phi(0)|\lambda_0\rangle \bigr|^2 .
  \label{eq:7.7}
\end{equation}
The spectral density $\rho(M^2)$ for a typical theory is plotted in
Fig.~\ref{fig:7.2}.  Note that the one-particle states contribute an isolated
delta function to the spectral density:
\begin{equation}
  \rho(M^2) = 2\pi\,\delta(M^2-m^2)\cdot Z
    + (\text{nothing else until } M^2 \ge (2m)^2),
  \label{eq:7.8}
\end{equation}
where $Z$ is some number given by the squared matrix element in
\eqref{eq:7.7}.  We refer to $Z$ as the \emph{field-strength renormalization}.
The quantity $m$ is the exact mass of a single particle---the exact energy
eigenvalue at rest.  This quantity will in general differ from the value of
the mass parameter that appears in the Lagrangian.  We will refer to the
parameter in the Lagrangian as $m_0$, the bare mass, and refer to $m$ as the
physical mass of the $\phi$ boson.  Only the physical mass $m$ is directly
observable.

The spectral decomposition \eqref{eq:7.6} yields the following form for the
Fourier transform of the two-point function:
\begin{equation}
  \begin{aligned}
    \int d^4x\ \langle\Omega|T\phi(x)\phi(0)|\Omega\rangle\ e^{ip\cdot x}
    &= \frac{iZ}{p^2-m^2+i\epsilon}
      + \int_{4m^2}^{\infty} \frac{dM^2}{2\pi}\ \rho(M^2)\
        \frac{i}{p^2-M^2+i\epsilon} .
  \end{aligned}
  \label{eq:7.9}
\end{equation}
The analytic structure of this function in the complex $p^2$-plane is shown
in Fig.~\ref{fig:7.3}.  The first term gives an isolated simple pole at
$p^2 = m^2$, while the second term gives a branch cut beginning at the
two-particle threshold $p^2 = 4m^2$.

\begin{figure}[t]
  \centering
  \includegraphics[width=0.5\textwidth]{fig7.3.png}
  \caption{Analytic structure in the complex $p^2$-plane of the Fourier
    transform of the two-point function for a typical theory.  The
    one-particle states contribute an isolated pole at the square of the
    particle mass.  States of two or more free particles give a branch cut,
    while bound states give additional poles.}
  \label{fig:7.3}
\end{figure}

For a field of spin $1/2$, the analog of \eqref{eq:7.9} is
\begin{equation}
  \int d^4x\ \langle\Omega|T\psi_\alpha(x)\bar{\psi}_\beta(0)|\Omega\rangle\
    e^{ip\cdot x}
  = \frac{i(\slashed{p}+m)\,Z_2}{p^2-m^2+i\epsilon}
    + \int_{4m^2}^{\infty} \frac{dM^2}{2\pi}\ \rho(M^2)\
      \frac{i(\slashed{p}+M)}{p^2-M^2+i\epsilon},
  \label{eq:7.12}
\end{equation}
where $Z_2$ is the field-strength renormalization for the Dirac field and
$\rho(M^2)$ is the corresponding spectral function.  For higher-spin fields,
the analog of \eqref{eq:7.9} is obtained by multiplying by the appropriate
polarization tensor.

\subsection{The Electron Self-Energy}

Let us now compute, in perturbation theory, the leading correction to the
electron propagator, and identify the field-strength renormalization constant
$Z_2$ and the mass shift $\delta m = m - m_0$.  We work in QED.  The lowest
order diagram contributing to the electron self-energy is
\begin{equation}
  \text{(one-loop electron self-energy diagram)} .
  \label{eq:7.15}
\end{equation}
Using the Feynman rules, its value is
\begin{equation}
  \begin{aligned}
    \Sigma_2(\slashed{p}) &= (-ie)^2\int \frac{d^4k}{(2\pi)^4}\
      \gamma^\mu\ \frac{i(\slashed{k}+m_0)}{k^2-m_0^2+i\epsilon}\
      \gamma^\nu\ \frac{-ig_{\mu\nu}}{(p-k)^2-\mu^2+i\epsilon}\\
    &= -e^2\int \frac{d^4k}{(2\pi)^4}\
      \frac{\gamma^\mu(\slashed{k}+m_0)\gamma_\mu}
        {[k^2-m_0^2+i\epsilon][(p-k)^2-\mu^2+i\epsilon]} .
  \end{aligned}
  \label{eq:7.16}
\end{equation}
Using the contraction identities
\begin{equation}
  \gamma^\mu\gamma_\mu = 4,\qquad
  \gamma^\mu\slashed{k}\gamma_\mu = -2\slashed{k},
\end{equation}
the numerator simplifies to
\begin{equation}
  \gamma^\mu(\slashed{k}+m_0)\gamma_\mu = 2m_0 - 2\slashed{k} .
\end{equation}
Introducing Feynman parameters to combine the denominators,
\begin{equation}
  \frac{1}{[k^2-m_0^2][(p-k)^2-\mu^2]}
  = \int_0^1 dx\ \frac{1}{[(k-xp)^2 - xp^2 + xp^2 + (1-x)m_0^2 + x\mu^2]^2}
  \cdots ,
\end{equation}
and shifting $\ell = k - xp$, we obtain
\begin{equation}
  \Sigma_2(\slashed{p}) = 2e^2\int_0^1 dx\ \int
    \frac{d^4\ell}{(2\pi)^4}\ \frac{-2x\slashed{p} + 4m_0}
      {[\ell^2 - \Delta + i\epsilon]^2},
  \label{eq:7.17}
\end{equation}
where $\Delta = -x(1-x)p^2 + x\mu^2 + (1-x)m_0^2$.  The integral over
$\ell$ is divergent.  To evaluate it, we first regulate it using the
Pauli--Villars procedure:
\begin{equation}
  \frac{1}{(p-k)^2-\mu^2+i\epsilon}
    \to \frac{1}{(p-k)^2-\mu^2+i\epsilon}
      - \frac{1}{(p-k)^2-\Lambda^2+i\epsilon} .
\end{equation}
The second term will have the same form as \eqref{eq:7.17}, but with $\mu$
replaced by $\Lambda$.  As in Section~6.3, we now Wick-rotate and substitute
the Euclidean variable $\ell_E^0 = -i\ell^0$.  This gives
\begin{equation}
  \int \frac{d^4\ell}{(2\pi)^4}\ \frac{1}{[\ell^2-\Delta]^2}
  = \frac{i}{(4\pi)^2}\int_0^{\infty} d\ell_E\
    \frac{\ell_E^3}{[\ell_E^2+\Delta]^2}
  = \frac{i}{(4\pi)^2}\ \log\Bigl( \frac{\Lambda^2}{\Delta} \Bigr),
  \label{eq:7.18}
\end{equation}
where $\Delta_\Lambda = -x(1-x)p^2 + x\Lambda^2 + (1-x)m_0^2 \to x\Lambda^2$
as $\Lambda\to\infty$.  The final result is therefore
\begin{equation}
  \Sigma_2(\slashed{p}) = \frac{\alpha}{2\pi}\int_0^1 dx\ (2m_0 - x\slashed{p})\
    \log\Bigl( \frac{x\Lambda^2}{(1-x)m_0^2 + x\mu^2 - x(1-x)p^2} \Bigr) .
  \label{eq:7.19}
\end{equation}
Before discussing the divergences in this expression, let us work out its
analytic behavior as a function of $p^2$.  The logarithm in \eqref{eq:7.19}
has a branch cut when its argument becomes negative, and for any fixed $x$
this will occur for sufficiently large $p^2$.  More exactly, the cut begins at
the point where
\begin{equation}
  (1-x)m_0^2 + x\mu^2 - x(1-x)p^2 = 0 .
\end{equation}
Solving this equation for $x$, we find the branch cut of $\Sigma_2(p^2)$
begins at $p^2 = (m_0+\mu)^2$, that is, at the threshold for creation of a
two-particle (electron plus photon) state.%
\footnote{In real QED, $\mu = 0$ and the two-particle branch cut merges with
the one-particle pole.  This subtlety plays a role in the full treatment of
the cancellation of infrared divergences, but it is beyond the scope of our
present analysis.}  The location of the branch cut is exactly where we would
expect from the K\"all\'en--Lehmann spectral representation.

We have now located the two-particle branch cut predicted by the
K\"all\'en--Lehmann representation, but we have not found the expected simple
pole at $p^2 = m^2$.  To find it we must actually include an infinite series
of Feynman diagrams.  Fortunately, this series will be easily summed.

Let us define a \emph{one-particle irreducible} (1PI) diagram to be any
diagram that cannot be split in two by removing a single line:
\begin{equation}
  \text{(a diagram with a self-energy insertion is 1PI)} .
\end{equation}
Let $-i\Sigma(\slashed{p})$ denote the sum of all 1PI diagrams with two
external fermion lines:
\begin{equation}
  -i\Sigma(\slashed{p}) = \text{(sum of all 1PI two-point diagrams)} .
  \label{eq:7.21}
\end{equation}
(Although each diagram has two external lines, the Feynman propagators for
these lines are not to be included in the expression for $\Sigma(p)$.)  To
leading order in $\alpha$ we see that $\Sigma = \Sigma_2$.

The Fourier transform of the two-point function can now be written as
\begin{equation}
  \int d^4x\ \langle\Omega|T\psi(x)\bar{\psi}(0)|\Omega\rangle\ e^{ip\cdot x}
  = \frac{i}{\slashed{p}-m_0}
    + \frac{i}{\slashed{p}-m_0}\ (-i\Sigma)\ \frac{i}{\slashed{p}-m_0}
    + \cdots .
  \label{eq:7.22}
\end{equation}
The first diagram has a simple pole at $p^2 = m_0^2$.  Each diagram in the
second class has a double pole at $p^2 = m_0^2$.  Each diagram in the third
class has a triple pole.  The behavior near $p^2 = m_0^2$ gets worse and worse
as we include more and more diagrams.  But fortunately, the sum of all the
diagrams forms a geometric series.  Note that $\Sigma(\slashed{p})$
commutes with $\slashed{p}$, since $\Sigma(\slashed{p})$ is a function only of
pure numbers and $\slashed{p}$.  In fact, we can consider
$\Sigma(\slashed{p})$ to be a function of $\slashed{p}$, writing
$p^2 = (\slashed{p})^2$.  Then we can rewrite each electron propagator as
$i/(\slashed{p}-m_0)$ and express the above series as
\begin{equation}
  \begin{aligned}
    \int d^4x\ \langle\Omega|T\psi(x)\bar{\psi}(0)|\Omega\rangle\ e^{ip\cdot x}
    &= \frac{i}{\slashed{p}-m_0}\ \Bigl[ 1 +
      \frac{\Sigma(\slashed{p})}{\slashed{p}-m_0}
      + \Bigl( \frac{\Sigma(\slashed{p})}{\slashed{p}-m_0} \Bigr)^2
      + \cdots \Bigr]\\
    &= \frac{i}{\slashed{p}-m_0-\Sigma(\slashed{p})} .
  \end{aligned}
  \label{eq:7.23}
\end{equation}
The full propagator has a simple pole, which is shifted away from $m_0$ by
$\Sigma(\slashed{p})$.  The location of this pole, the physical mass $m$, is
the solution of the equation
\begin{equation}
  \bigl[ \slashed{p} - m_0 - \Sigma(\slashed{p}) \bigr]
    \Big|_{\slashed{p}=m} = 0 .
  \label{eq:7.24}
\end{equation}
Notice that, if $\Sigma(\slashed{p})$ is defined by the convention
\eqref{eq:7.21}, then a positive contribution to $\Sigma$ yields a positive
shift of the electron mass.  Close to the pole, the denominator of
\eqref{eq:7.23} has the form
\begin{equation}
  (\slashed{p}-m)\cdot \Bigl( 1 - \frac{d\Sigma}{d\slashed{p}}
    \Big|_{\slashed{p}=m} \Bigr) + \mathcal{O}\bigl( (\slashed{p}-m)^2 \bigr)
  = (\slashed{p}-m)\cdot Z_2^{-1} + \cdots .
  \label{eq:7.25}
\end{equation}
Thus the full electron propagator has a single-particle pole of just the form
\eqref{eq:7.12}, with $m$ given by \eqref{eq:7.24} and
\begin{equation}
  Z_2^{-1} = 1 - \frac{d\Sigma}{d\slashed{p}}\Big|_{\slashed{p}=m} .
  \label{eq:7.26}
\end{equation}
Our explicit calculation of $\Sigma_2$ allows us to compute the first
corrections to $m$ and $Z_2$.  Let us begin with $m$.  To order $\alpha$, the
mass shift is
\begin{equation}
  \delta m = m - m_0 = \Sigma_2(\slashed{p}=m) \approx \Sigma_2(\slashed{p}=m_0) .
  \label{eq:7.27}
\end{equation}
Carrying out the integral, one finds that $\delta m$ is logarithmically
divergent (ultraviolet) and also contains an infrared logarithm.  The
ultraviolet divergence is absorbed by a renormalization of the bare mass,
while the infrared divergence cancels against bremsstrahlung, as we discussed
in Chapter~6.  Similarly, the derivative of $\Sigma_2$ at
$\slashed{p} = m$ determines $Z_2$; to order $\alpha$,
\begin{equation}
  Z_2 = 1 - \frac{d\Sigma_2}{d\slashed{p}}\Big|_{\slashed{p}=m}
    + \mathcal{O}(\alpha^2) .
  \label{eq:7.29}
\end{equation}

\section{The LSZ Reduction Formula}

We would now like to express the $S$-matrix elements of our theory---for
example, the scattering amplitude for a $2\to2$ process---in terms of
Fourier transforms of vacuum expectation values of time-ordered products of
fields, that is, correlation functions.  This relation is provided by the LSZ
(Lehmann--Symanzik--Zimmermann) reduction formula, which we derive in this
section.%

\subsection*{Derivation of the LSZ Formula}

Consider the amplitude for a single-particle state of definite momentum
$\mathbf{p}$ to evolve into the same state at a later time.  In the
interacting theory, the state $|\mathbf{p}\rangle$ that evolves freely in the
distant past is not the same as the eigenstate of the full Hamiltonian with
definite momentum.  Rather, the states that appear in the spectral
decomposition are the physical one-particle states, with the normalization
\begin{equation}
  \langle\mathbf{p}|\mathbf{q}\rangle = 2E_{\mathbf{p}}\,(2\pi)^3\,
    \delta^{(3)}(\mathbf{p}-\mathbf{q}) .
\end{equation}
The field $\phi(x)$ has a nonzero matrix element between the vacuum and this
physical one-particle state,
\begin{equation}
  \langle\Omega|\phi(x)|\mathbf{p}\rangle = \sqrt{Z}\,
    e^{-ip\cdot x},
  \label{eq:7.35}
\end{equation}
where $Z$ is the field-strength renormalization constant of
Eq.~\eqref{eq:7.8}.

To derive the reduction formula, consider the correlation function with two
external particles:
\begin{equation}
  \langle\Omega|T\{\phi(x_1)\phi(x_2)\cdots\}|\Omega\rangle .
  \label{eq:7.41}
\end{equation}
Insert a complete set of states between the fields, and analyze the time
integrals by separating the coordinates into regions where the particles are
well separated (the far past and far future) and a bounded region where they
interact.  The result of this analysis is that the leading singularity in the
energy variables, when the wavepackets are taken to definite momenta, gives
the $S$-matrix element.  By interchanging the order of the limits, one obtains
the LSZ reduction formula:
\begin{equation}
  \begin{aligned}
    &\Bigl[ \prod_{i=1}^{n} \int d^4x_i\ e^{ip_i\cdot x_i}\
      \frac{i}{p_i^2-m^2+i\epsilon}\ \sqrt{Z}\ \Bigr]\
    \Bigl[ \prod_{j=1}^{m} \int d^4y_j\ e^{-ik_j\cdot y_j}\
      \frac{i}{k_j^2-m^2+i\epsilon}\ \sqrt{Z}\ \Bigr]\\
    &\quad \times \langle\Omega|T\{\phi(x_1)\cdots\phi(x_n)\phi(y_1)\cdots
      \phi(y_m)\}|\Omega\rangle\\
    &\xrightarrow{\text{each } p_i^0\to +E_{\mathbf{p}_i};\ \text{each }
      k_j^0\to +E_{\mathbf{k}_j}}
    \langle\mathbf{p}_1\cdots\mathbf{p}_n|S|\mathbf{k}_1\cdots
      \mathbf{k}_m\rangle .
  \end{aligned}
  \label{eq:7.42}
\end{equation}
The quantity $Z$ that appears in this equation is exactly the
field-strength renormalization constant, defined in Section~7.1 as the
residue of the single-particle pole in the two-point function of fields.  Each
distinct particle will have a distinct renormalization factor $Z$, obtained
from its own two-point function.  For higher-spin fields, each factor of
$\sqrt{Z}$ comes with a polarization factor such as $u^s(p)$, as in
Eq.~\eqref{eq:7.12}.  The polarization $s$ must be summed over in the second
line of \eqref{eq:7.42}.

In words, the LSZ formula says that an $S$-matrix element can be computed as
follows.  Compute the appropriate Fourier-transformed correlation function,
look at the region of momentum space where the external particles are near
the mass shell, and identify the coefficient of the multiparticle pole.  For
fields with spin, one must also multiply by a polarization spinor (like
$u^s(p)$) or vector (like $\epsilon^r(k)$) to project out the desired spin
state.

Our next goal is to express this procedure in the language of Feynman
diagrams.  Let us analyze the relation between the diagrammatic expansion of
the scalar field four-point function and the $S$-matrix element for
2-particle $\to$ 2-particle scattering.  We will consider explicitly the fully
connected Feynman diagrams contributing to the correlator.  By a similar
analysis, it is easy to confirm that disconnected diagrams should be
disregarded because they do not have the singularity structure, with a product
of four poles, indicated on the right-hand side of \eqref{eq:7.42}.

The exact four-point function
\begin{equation}
  \Bigl( \prod_{i=1}^{2}\int d^4x_i\ e^{ip_i\cdot x_i}\ \Bigr)\
    \Bigl( \prod_{j=1}^{2}\int d^4y_j\ e^{-ik_j\cdot y_j}\ \Bigr)\
    \langle\Omega|T\{\phi(x_1)\phi(x_2)\phi(y_1)\phi(y_2)\}|\Omega\rangle
\end{equation}
has the general form shown in Fig.~\ref{fig:7.4}.  In this figure, we have
indicated explicitly the diagrammatic corrections on each leg; the shaded
circle in the center represents the sum of all amputated four-point diagrams.

\begin{figure}[t]
  \centering
  \includegraphics[width=0.55\textwidth]{fig7.4.png}
  \caption{Structure of the exact four-point function in scalar field theory.
    The external legs carry the field-strength renormalization factors
    $\sqrt{Z}$, and the shaded circle is the sum of all amputated 1PI
    four-point diagrams.}
  \label{fig:7.4}
\end{figure}

We can sum up the corrections to each external leg just as we did for the
electron propagator in the previous section.  Let $-iM^2(p^2)$ denote the sum
of all one-particle-irreducible (1PI) insertions into the scalar propagator:
\begin{equation}
  -iM^2(p^2) = \text{(sum of 1PI self-energy insertions)} .
\end{equation}
Then the exact propagator can be written as a geometric series and summed as
in Eq.~\eqref{eq:7.23}:
\begin{equation}
  \frac{i}{p^2-m_0^2-M^2(p^2)} .
  \label{eq:7.43}
\end{equation}
Notice that, as in the case of the electron propagator, our sign convention
for the 1PI self-energy $M^2(p^2)$ implies that a positive contribution to
$M^2(p^2)$ corresponds to a positive shift of the scalar particle mass.  If we
expand each resummed propagator about the physical particle pole, we see that
each external leg of the four-point amplitude contributes
\begin{equation}
  \frac{iZ}{p^2-m^2-M^2} \xrightarrow{p^0\to E_{\mathbf{p}}}
    \frac{iZ}{p^2-m^2} + (\text{regular}) .
  \label{eq:7.44}
\end{equation}
Thus, the sum of diagrams contains a product of four poles:
\begin{equation}
  \frac{iZ}{p_1^2-m^2}\ \frac{iZ}{p_2^2-m^2}\ \frac{iZ}{k_1^2-m^2}\
    \frac{iZ}{k_2^2-m^2}\ \times\ \text{(amputated, connected amplitude)} .
\end{equation}
The LSZ formula \eqref{eq:7.42} instructs us to multiply the correlator by
$\prod (i\sqrt{Z}/(p^2-m^2+i\epsilon))$ and then take the residues at the
poles.  The factors of $i/(p^2-m^2)$ cancel against the propagators on the
external legs, and we are left with
\begin{equation}
  \langle\mathbf{p}_1\mathbf{p}_2|S|\mathbf{k}_1\mathbf{k}_2\rangle
  = \text{sum of all amputated, connected diagrams},
  \label{eq:7.45}
\end{equation}
which is the master formula \eqref{eq:4.103}.  This completes the derivation
of the Feynman rules for $S$-matrix elements.

\section{The Optical Theorem}

The $S$-matrix satisfies a number of general identities that follow from
unitarity.  The most important of these is the optical theorem, which relates
the imaginary part of the forward scattering amplitude to the total cross
section.  In this section we will prove this theorem and show how it arises in
the Feynman diagram expansion.

\subsection{Unitarity and the Optical Theorem}

The $S$-matrix is unitary: $S^{\dagger}S = 1$.  Writing $S = 1+iT$, this
implies
\begin{equation}
  -i(T - T^{\dagger}) = T^{\dagger}T .
\end{equation}
Taking matrix elements between an initial state $|i\rangle$ and a final state
$|f\rangle$, and inserting a complete set of states on the right, we obtain
\begin{equation}
  -i\bigl( \langle f|T|i\rangle - \langle f|T^{\dagger}|i\rangle \bigr)
  = \sum_n \langle f|T^{\dagger}|n\rangle\langle n|T|i\rangle .
\end{equation}
For the forward scattering amplitude, $|f\rangle = |i\rangle$, this gives
\begin{equation}
  2\,\mathrm{Im}\,\mathcal{M}(i\to i)
  = \sum_n \int d\Pi_n\ |\mathcal{M}(i\to n)|^2 .
  \label{eq:7.49}
\end{equation}
Using the relation between $|\mathcal{M}|^2$ and the cross section, we obtain
the optical theorem:
\begin{equation}
  \mathrm{Im}\,\mathcal{M}(i\to i) = 2E_{\mathrm{cm}}\,p_{\mathrm{cm}}\,
    \sigma_{\mathrm{total}}(i\to\text{anything}),
  \label{eq:7.50}
\end{equation}
where $E_{\mathrm{cm}}$ is the total center-of-mass energy and $p_{\mathrm{cm}}$
is the momentum of either particle in the center-of-mass frame.  This equation
relates the forward scattering amplitude to the total cross section for
production of all final states.  Since the imaginary part of the forward
scattering amplitude gives the attenuation of the forward-going wave as the
beam passes through the target, it is natural that this quantity should be
proportional to the probability of scattering.  Identity \eqref{eq:7.50}
gives the precise connection.

\subsection{The Optical Theorem for Feynman Diagrams}

Let us now investigate how this identity for the imaginary part of an
$S$-matrix element arises in the Feynman diagram expansion.  It is easily
checked (in QED, for example) that each diagram contributing to an $S$-matrix
element $\mathcal{M}$ is purely real unless some denominators vanish, so that
the $i\epsilon$ prescription for treating the poles becomes relevant.  A
Feynman diagram thus yields an imaginary part for $\mathcal{M}$ only when the
virtual particles in the diagram go on-shell.  We will now show how to isolate
and compute this imaginary part.

For our present purposes, let us define $\mathcal{M}$ by the Feynman rules
for perturbation theory.  This allows us to consider $\mathcal{M}(s)$ as an
analytic function of the complex variable $s = E^2_{\mathrm{cm}}$, even though
$S$-matrix elements are defined only for external particles with real momenta.

We first demonstrate that the appearance of an imaginary part of
$\mathcal{M}(s)$ always requires a branch cut singularity.  Let $s_0$ be the
threshold energy for production of the lightest multiparticle state.  For real
$s$ below $s_0$ the intermediate state cannot go on-shell, so
$\mathcal{M}(s)$ is real.  Thus, for real $s < s_0$, we have the identity
\begin{equation}
  \mathcal{M}(s) = \bigl[ \mathcal{M}(s) \bigr]^{*} .
  \label{eq:7.51}
\end{equation}
Each side of this equation is an analytic function of $s$, so it can be
analytically continued to the entire complex $s$ plane.  In particular, near
the real axis for $s > s_0$, Eq.~\eqref{eq:7.51} implies
\begin{equation}
  \mathrm{Re}\,\mathcal{M}(s+i\epsilon) = \mathrm{Re}\,\mathcal{M}(s-i\epsilon);
  \qquad
  \mathrm{Im}\,\mathcal{M}(s+i\epsilon) = -\mathrm{Im}\,\mathcal{M}(s-i\epsilon).
\end{equation}
There is a branch cut across the real axis, starting at the threshold energy
$s_0$; the discontinuity across the cut is
\begin{equation}
  \mathrm{Disc}\,\mathcal{M}(s) = 2i\,\mathrm{Im}\,\mathcal{M}(s+i\epsilon) .
\end{equation}
Usually it is easier to compute the discontinuity of a diagram than to
compute the imaginary part directly.  The $i\epsilon$ prescription in the
Feynman propagator indicates that physical scattering amplitudes should be
evaluated above the cut, at $s+i\epsilon$.

We already saw in Section~7.1 that the electron self-energy diagram has a
branch cut beginning at the physical electron-photon threshold.  Let us now
study more general one-loop diagrams, and show that their discontinuities give
precisely the imaginary parts required by \eqref{eq:7.49}.  The generalization
of this result to multiloop diagrams has been proven by Cutkosky,%
\footnote{R.~E.~Cutkosky, \emph{J.\ Math.\ Phys.} \textbf{1}, 429 (1960).} who
showed in the process that the discontinuity of a Feynman diagram across its
branch cut is always given by a simple set of cutting rules.%
\footnote{These rules are simple only for singularities in the physical
region.  Away from the physical region, the singularities of three- and
higher-point amplitudes can become quite intricate.  This subject is reviewed
in R.~J.~Eden, P.~V.~Landshoff, D.~I.~Olive, and J.~C.~Polkinghorne,
\emph{The Analytic S-Matrix} (Cambridge University Press, 1966).}

We begin by checking \eqref{eq:7.49} in $\phi^4$ theory.  Since the
right-hand side of \eqref{eq:7.49} begins in order $\lambda^2$, we expect that
$\mathrm{Im}\,\mathcal{M}$ should also receive its first contribution from
higher-order diagrams.  Consider, then, the order-$\lambda^2$ diagram
\begin{equation}
  \delta\mathcal{M}(s) = \frac{(-i\lambda)^2}{2}\ \int
    \frac{d^4q}{(2\pi)^4}\ \frac{i}{(k/2-q)^2-m^2+i\epsilon}\
    \frac{i}{(k/2+q)^2-m^2+i\epsilon} ,
  \label{eq:7.52}
\end{equation}
with a loop in the $s$-channel ($k = k_1+k_2$ is the total momentum).  (It is
easy to check that the corresponding $t$- and $u$-channel diagrams have no
branch cut singularities for $s$ above threshold.)  When this integral is
evaluated using the methods of Section~6.3, the Wick rotation produces an
extra factor of $i$, so that, below threshold, $\delta\mathcal{M}$ is purely
real.

We would like to verify that the integral \eqref{eq:7.52} has a discontinuity
across the real axis in the physical region $k^0 > 2m$.  It is easiest to
identify this discontinuity by computing the integral for $k^0 < 2m$, then
increasing $k^0$ by analytic continuation.  Let us work in the center-of-mass
frame, where $k = (k^0, \mathbf{0})$.  Then the integrand of \eqref{eq:7.52}
has four poles in the integration variable $q^0$, at the locations
\begin{equation}
  q^0 = \frac{1}{2}k^0 \pm (E_{\mathbf{q}} - i\epsilon),
  \qquad
  q^0 = -\frac{1}{2}k^0 \pm (E_{\mathbf{q}} - i\epsilon) .
\end{equation}
Two of these poles lie above the real $q^0$ axis and two lie below.  We will
close the integration contour downward and pick up the residues of the poles
in the lower half-plane.  Of these, only the pole at
$q^0 = -\tfrac{1}{2}k^0 + E_{\mathbf{q}}$ will contribute to the
discontinuity.  Note that picking up the residue of this pole is equivalent
to replacing
\begin{equation}
  \frac{1}{(k/2+q)^2-m^2+i\epsilon}
    \to -2\pi i\,\delta\bigl( (k/2+q)^2 - m^2 \bigr)
  \label{eq:7.53}
\end{equation}
under the $dq^0$ integral.

The contribution of this pole yields the integral
\begin{equation}
  \begin{aligned}
    i\delta\mathcal{M} &= -2\pi i\ \frac{\lambda^2}{2}\ \int
      \frac{d^3q}{(2\pi)^4}\ \frac{1}{2E_{\mathbf{q}}}\
      \frac{i}{(k^0-E_{\mathbf{q}})^2}\\
    &= -\frac{\lambda^2}{2}\ \frac{4\pi}{(2\pi)^4}\ \int_0^{\infty}
      dE_{\mathbf{q}}\ E_{\mathbf{q}}|\mathbf{q}|\ \frac{1}{2E_{\mathbf{q}}\,k^0\,
        (k^0-2E_{\mathbf{q}})} .
  \end{aligned}
  \label{eq:7.54}
\end{equation}
The integrand in the second line has a pole at $E_{\mathbf{q}} = k^0/2$.
When $k^0 < 2m$, this pole does not lie on the integration contour, so
$\delta\mathcal{M}$ is manifestly real.  When $k^0 > 2m$, however, the pole
lies just above or below the contour of integration, depending upon whether
$k^0$ is given a small positive or negative imaginary part.  Thus the integral
acquires a discontinuity between $k^2 + i\epsilon$ and $k^2 - i\epsilon$.  To
compute this discontinuity, apply
\begin{equation}
  \frac{1}{k^0 - 2E_{\mathbf{q}} \pm i\epsilon}
  = \mathrm{P}\frac{1}{k^0 - 2E_{\mathbf{q}}}
    \mp i\pi\,\delta(k^0 - 2E_{\mathbf{q}})
\end{equation}
(where P denotes the principal value).  The discontinuity is given by
replacing the pole with a delta function.  This in turn is equivalent to
replacing the original propagator by a delta function:
\begin{equation}
  \frac{1}{(k/2-q)^2-m^2+i\epsilon}
    \to -2\pi i\,\delta\bigl( (k/2-q)^2 - m^2 \bigr) .
  \label{eq:7.55}
\end{equation}
Let us now retrace our steps and see what we have proved.  Go back to the
original integral \eqref{eq:7.52}, relabel the momenta on the two propagators
as $p_1$, $p_2$ and substitute
\begin{equation}
  \frac{1}{p_i^2 - m^2 + i\epsilon}
    \to -2\pi i\,\delta(p_i^2 - m^2) .
  \label{eq:7.56}
\end{equation}
We have shown that the discontinuity of the integral is computed by replacing
each of the two propagators by a delta function.  The discontinuity of
$\mathcal{M}$ comes only from the region of the $d^4p_i$ integral in which the
two delta functions are simultaneously satisfied.  By integrating over the
delta functions, we put the momenta $p_i$ on shell and convert the integrals
$d^4p_i$ into integrals over relativistic phase space.  What is left over in
expression \eqref{eq:7.52} is just the factor $\lambda^2$, the square of the
leading-order scattering amplitude, and the symmetry factor $(1/2)$, which
can be reinterpreted as the appropriate phase-space normalization.  This is
precisely the optical theorem \eqref{eq:7.49} for the two-body intermediate
state.  This establishes the optical theorem for one-loop diagrams.

\section{The Ward-Takahashi Identity}

In this section we will give a formal proof of the Ward identity
\eqref{eq:5.79}, and in fact of a more general relation, the
Ward--Takahashi identity.  Consider a QED amplitude $\mathcal{M}$ with an
external photon of momentum $k$ and polarization $\epsilon(k)$.  Writing
$\mathcal{M} = \epsilon^{*\mu}(k)\mathcal{M}_\mu(k)$, we want to show that
\begin{equation}
  k^\mu\mathcal{M}_\mu(k) = 0 .
  \label{eq:7.64}
\end{equation}
Our proof proceeds diagram by diagram.  Consider first a through-going
electron line with $n$ vertices, to which an additional external photon is to
be attached.  The electron propagators have momenta $p, p_1 = p+q_1,
p_2 = p_1+q_2$, and so on up to $p_f = p_{n-1}+q_n$.  If there are $n$
vertices, we can insert our photon in $n+1$ different places.  Suppose we
insert it after the $i$th vertex.

Let us now look at the values of these diagrams, with the polarization vector
$\epsilon^\mu(k)$ replaced by $k^\mu$.  The product of $k^\mu$ with the new
vertex is conveniently written:
\begin{equation}
  -ie\,k^\mu\gamma_\mu = -ie\bigl[ (\slashed{p}_i+\slashed{k}-m)
    - (\slashed{p}_i-m) \bigr] .
  \label{eq:7.65}
\end{equation}
Multiplying by the adjacent electron propagators, we obtain
\begin{equation}
  \frac{i}{\slashed{p}_i+\slashed{k}-m}\ (-ie)\ \frac{i}{\slashed{p}_i-m}
  = ie\Bigl( \frac{1}{\slashed{p}_i-m}
    - \frac{1}{\slashed{p}_i+\slashed{k}-m} \Bigr) .
\end{equation}
The diagram with the photon inserted at position $i$ therefore has a structure
in which the first term of this expression cancels the second term of the
diagram with the photon inserted at position $i-1$.  A similar cancellation
occurs between any other pair of diagrams with adjacent insertions.  When we
sum over all possible insertion points along the line, everything cancels
except the unpaired terms at the ends.  The unpaired term coming from
insertion after the last vertex (on the far left) is just $e$ times the value
of the original diagram; the other unpaired term, from insertion before the
first vertex, is identical except for a minus sign and the replacement of $p$
by $p+k$ everywhere.  Diagrammatically, our result is
\begin{equation}
  \sum_{\text{insertion points}} \text{(diagram with photon inserted)}
  = e\Bigl[ \text{(original diagram at momentum }p\text{)}
    - \text{(original diagram at momentum }p+k\text{)} \Bigr].
  \label{eq:7.66}
\end{equation}
In each diagram on the left-hand side of \eqref{eq:7.66}, the momentum
entering the electron line is $p$ and the momentum exiting is $q$.  According
to the LSZ formula, we can extract from each diagram a contribution to an
$S$-matrix element by taking the coefficient of the product of poles
\begin{equation}
  \frac{1}{\slashed{p}-m}\ \frac{1}{\slashed{q}-m} .
\end{equation}
The terms on the right-hand side of \eqref{eq:7.66} each contain one of these
poles, but neither contains both poles.  Thus the right-hand side of
\eqref{eq:7.66} contributes nothing to the $S$-matrix.%
\footnote{This step of the argument is straightforward only if we have
arranged the perturbation series so that the propagator contains $m$ rather
than $m_0$.  We will do this in Chapter~10.}

To complete the proof of the Ward--Takahashi identity, we must consider the
case in which our photon attaches to an internal electron loop.  Before the
insertion of the photon, a typical loop has $n$ vertices and $n$ propagators
with momenta $p_1, p_1+q_2 = p_2$, and so on up to $p_n$.  Suppose now that we
insert the photon $\gamma(k)$ between vertices $i$ and $i+1$.  We now have an
additional momentum $k$ running around the loop from the new vertex.  To
evaluate the sum over all such insertions into the loop, apply identity
\eqref{eq:7.65} to each diagram.  The first term of one diagram will be
canceled by one of the terms from the diagram with the photon inserted at the
next vertex.  Similar cancellations take place between terms from other pairs
of adjacent insertions.  When we sum over all $n$ insertion points we are left
with
\begin{equation}
  \text{(diagram with photon inserted between vertex }n\text{ and vertex }
    1\text{)}
  + \text{(diagram with photon inserted between vertex }1\text{ and vertex }
    2\text{)},
  \label{eq:7.67}
\end{equation}
which differ only by a shift of the loop momentum by $k$.  Shifting the
integration variable from $p_1$ to $p_1+k$ in the second term, we see that the
two remaining terms cancel.  Thus the diagrams in which the photon is inserted
along a closed loop add up to zero.

We are now ready to assemble the pieces of the proof.  Suppose that the
amplitude $\mathcal{M}$ has $2n$ external electron lines, $n$ incoming and $n$
outgoing.  Label the incoming momenta $p_i$ and the outgoing momenta
$q_i$:
\begin{equation}
  \mathcal{M}_\mu(k;\ p_1\cdots p_n;\ q_1\cdots q_n) .
\end{equation}
(The amplitude can also involve an arbitrary number of additional external
photons.)  The amplitude $\mathcal{M}_0$ lacks the photon $\gamma(k)$ but is
otherwise identical.  To form $k^\mu\mathcal{M}_\mu$ from $\mathcal{M}_0$ we
must sum over all diagrams that contribute to $\mathcal{M}_0$, and for each
diagram, sum over all points at which the photon could be inserted.  Summing
over insertion points along an internal loop in any diagram gives zero.
Summing over insertion points along a through-going line in any diagram gives
a contribution of the form \eqref{eq:7.66}.  Summing over all insertion
points for any particular diagram, we obtain
\begin{equation}
  \text{(difference of diagrams with one external momentum shifted)} .
\end{equation}
Summing over all such diagrams, we finally obtain
\begin{equation}
  \begin{aligned}
    k_\mu\mathcal{M}^{\mu}(k;\ p_1\cdots p_n;\ q_1\cdots q_n)
    &= e\ \sum_{i} \Bigl[ \mathcal{M}_0(p_1\cdots p_n;\ q_1\cdots
      (q_i-k)\cdots)\\
    &\quad - \mathcal{M}_0((p_i+k)\cdots p_n;\ q_1\cdots q_n) \Bigr].
  \end{aligned}
  \label{eq:7.68}
\end{equation}
This is the Ward--Takahashi identity for correlation functions in QED.  We
saw below \eqref{eq:7.66} that the right-hand side does not contribute to the
$S$-matrix; thus in the special case where $\mathcal{M}$ is an $S$-matrix
element, Eq.~\eqref{eq:7.68} reduces to the Ward identity \eqref{eq:7.64}.

Before discussing this identity further, we should mention a potential flaw in
the above proof.  In order to find the necessary cancellation in Eq.
\eqref{eq:7.67}, we had to shift the integration variable by a constant.  If
the integral diverges, however, this shift is not permissible.  Similarly,
there may be divergent loop-momentum integrals in the expressions leading to
Eq.~\eqref{eq:7.66}.  Here there is no explicit shift in the proof, but in
practice one does generally perform a shift while evaluating the integrals.
In either case, ultraviolet divergences can potentially invalidate the
Ward--Takahashi identity.  We will see an example of this problem, as well as
a general solution to it, in the next section.

\section{Renormalization of the Electric Charge}

The last topic of this chapter is the computation of the vacuum polarization
diagram, the fourth diagram of \eqref{eq:6.1}.  This diagram is also known as
the correction to the photon propagator.  Let
$-i\Pi^{\mu\nu}(q)$ denote the sum of all 1PI insertions into the photon
propagator:
\begin{equation}
  -i\Pi^{\mu\nu}(q) = \text{(sum of all 1PI photon self-energy diagrams)} .
  \label{eq:7.71}
\end{equation}
Lorentz invariance and current conservation constrain the form of
$\Pi^{\mu\nu}(q)$:
\begin{equation}
  \Pi^{\mu\nu}(q) = (q^2g^{\mu\nu} - q^\mu q^\nu)\,\Pi(q^2) .
  \label{eq:7.72}
\end{equation}
(Any term proportional to $g^{\mu\nu}$ without the compensating $q^\mu q^\nu$
would violate the Ward identity and give the photon a mass.)  The exact photon
propagator is obtained by summing the geometric series of 1PI insertions:
\begin{equation}
  \frac{-ig_{\mu\nu}}{q^2}\ \Bigl[ 1 + \Pi(q^2) + \Pi(q^2)^2 + \cdots \Bigr]
  = \frac{-ig_{\mu\nu}}{q^2\,[1 - \Pi(q^2)]} .
  \label{eq:7.73}
\end{equation}
In any $S$-matrix element calculation, at least one end of this exact
propagator will connect to a fermion line.  When we sum over all places along
the line where it could connect, we must find, according to the Ward identity,
that terms proportional to $q_\mu$ or $q_\nu$ vanish.  For the purposes of
computing $S$-matrix elements, therefore, we can abbreviate
\begin{equation}
  \frac{-ig_{\mu\nu}}{q^2\,[1-\Pi(q^2)]}
  \xrightarrow{\text{Ward identity}}
  \frac{-ig_{\mu\nu}}{q^2}\ \frac{1}{1-\Pi(q^2)} .
  \label{eq:7.75}
\end{equation}
Notice that as long as $\Pi(q^2)$ is regular at $q^2 = 0$, the exact
propagator always has a pole at $q^2 = 0$.  In other words, the photon remains
absolutely massless at all orders in perturbation theory.  This claim depends
strongly on our use of the Ward identity in \eqref{eq:7.73}.  If, for example,
$\Pi^{\mu\nu}(q)$ contained a term $M^2g^{\mu\nu}$ (with no compensating
$q^\mu q^\nu$ term), the photon mass would be shifted to $M$.  The residue of
the $q^2 = 0$ pole is
\begin{equation}
  \frac{1}{1 - \Pi(0)} = Z_3 .
  \label{eq:7.75b}
\end{equation}
The amplitude for any low-$q^2$ scattering process will be shifted by this
factor, relative to the tree-level approximation:
\begin{equation}
  e_0^2 \to e_0^2\,Z_3 .
\end{equation}
Since a factor of $e$ lies at each end of the photon propagator, we can
conveniently account for this shift by making the replacement
$e_0 \to \sqrt{Z_3}\,e_0$.  This replacement is called \emph{charge
renormalization}; it is in many ways analogous to the mass renormalization
introduced in Section~7.1.  Note in particular that the ``physical'' electron
charge measured in experiments is $\sqrt{Z_3}\,e_0$.  We will therefore shift
our notation and call this quantity simply $e$.  From now on we will refer to
the ``bare'' charge (the quantity that multiplies $\gamma^\mu$ in the
Lagrangian) as $e_0$.  We then have
\begin{equation}
  (\text{physical charge}) = e = \sqrt{Z_3}\,e_0 = \sqrt{Z_3}\,
    (\text{bare charge}) .
  \label{eq:7.76}
\end{equation}
To lowest order, $Z_3 = 1$ and $e = e_0$.

In addition to this constant shift in the strength of the electric charge,
$\Pi(q^2)$ has another effect.  Consider a scattering process with nonzero
$q^2$, and suppose that we have computed $\Pi(q^2)$ to leading order in
$\alpha$: $\Pi(q^2) \approx \Pi_2(q^2)$.  The amplitude for the process will
then involve the quantity
\begin{equation}
  \frac{e_0^2}{q^2}\ \frac{1}{1-\Pi(q^2)} = \frac{e^2}{q^2}\
    \frac{1}{1 - [\Pi_2(q^2) - \Pi_2(0)]} .
\end{equation}
(Swapping $e^2$ for $e_0^2$ does not matter to lowest order.)  The quantity in
parentheses can be interpreted as a $q^2$-dependent electric charge.  The
full effect of replacing the tree-level photon propagator with the exact
photon propagator is therefore to replace
\begin{equation}
  \alpha = \frac{e_0^2}{4\pi}
    \to \alpha_{\mathrm{eff}}(q^2) = \frac{e^2}{4\pi}\
      \frac{1}{1 - [\Pi_2(q^2)-\Pi_2(0)]}
    = \frac{\alpha}{1 - [\Pi_2(q^2)-\Pi_2(0)]} .
  \label{eq:7.77}
\end{equation}
(To leading order we could just as well bring the $\Pi$-terms into the
numerator; but we will see in Chapter~12 that in this form, the expression is
true to all orders when $\Pi_2$ is replaced by $\Pi$.)

\subsection{Computation of $\Pi_2$}

Having worked so hard to interpret $\Pi_2(q^2)$, we had better calculate it.
Going back to \eqref{eq:7.71}, we have
\begin{equation}
  \begin{aligned}
    -i\Pi^{\mu\nu}_2(q) &= -(-ie)^2 \int \frac{d^4k}{(2\pi)^4}\
      \mathrm{tr}\Bigl[ \gamma^\mu\ \frac{i(\slashed{k}+m)}{k^2-m^2}\
        \gamma^\nu\ \frac{i(\slashed{k}+\slashed{q}+m)}
          {(k+q)^2-m^2} \Bigr]\\
    &= -4e^2 \int \frac{d^4k}{(2\pi)^4}\
      \frac{k^\mu(k+q)^\nu + k^\nu(k+q)^\mu - g^{\mu\nu}
        (k\cdot(k+q) - m^2)}
        {(k^2-m^2)((k+q)^2-m^2)} .
  \end{aligned}
  \label{eq:7.78}
\end{equation}
We have written $e$ and $m$ instead of $e_0$ and $m_0$ for convenience, since
the difference would give only an order-$\alpha^2$ contribution to
$\Pi_2$.

Now introduce a Feynman parameter to combine the denominator factors:
\begin{equation}
  \frac{1}{(k^2-m^2)((k+q)^2-m^2)}
  = \int_0^1 dx\ \frac{1}{(k^2 + 2xk\cdot q + xq^2 - m^2)^2}
  = \int_0^1 dx\ \frac{1}{(\ell^2 + x(1-x)q^2 - m^2)^2},
\end{equation}
where $\ell = k + xq$.  In terms of $\ell$, the numerator of \eqref{eq:7.78}
is
\begin{equation}
  \text{Numerator} = 2\ell^\mu\ell^\nu - 2x(1-x)q^\mu q^\nu
    + g^{\mu\nu}\bigl( m^2 + x(1-x)q^2 \bigr)
    + (\text{terms linear in } \ell) .
\end{equation}
Performing a Wick rotation and substituting $\ell_E^0 = -i\ell^0$, we obtain
\begin{equation}
  -i\Pi^{\mu\nu}_2(q) = 4e^2 \int_0^1 dx\ \int
    \frac{d^4\ell_E}{(2\pi)^4}\ \frac{1}{(\ell_E^2+\Delta)^2}\
    \Bigl[ 2\ell_E^\mu\ell_E^\nu - g^{\mu\nu}\ell_E^2 + \cdots \Bigr],
  \label{eq:7.79}
\end{equation}
where $\Delta = m^2 - x(1-x)q^2$.  This integral is very badly ultraviolet
divergent.  If we were to cut it off at $\ell_E = \Lambda$, we would find for
the leading term,
\begin{equation}
  i\Pi^{\mu\nu}_2(q) \propto e^2\Lambda^2 g^{\mu\nu},
\end{equation}
with no compensating $q^\mu q^\nu$ term.  This result violates the Ward
identity; it would give the photon an infinite mass $M \propto e\Lambda$.

Our proof of the Ward identity has failed, in precisely the way anticipated at
the end of the previous section: The shift of the integration variable in
\eqref{eq:7.67} is not permissible when the integral is divergent.  In our
present calculation, the failure of the Ward identity is catastrophic: It
leads to an infinite photon mass,%
\footnote{We could still make the observed photon mass zero by adding a
compensating infinite photon mass term to the Lagrangian.  More generally, we
could add terms to the Lagrangian to make the Ward identity valid for any
$n$-point correlation function.  This procedure would give the same results as
the one we are about to follow, but would be much more complicated.} in
conflict with experiment.  Fortunately, there is a way to rescue the Ward
identity.

In the above analysis we regulated the divergent integral in the most
straightforward and most naive way: by cutting it off at a large momentum
$\Lambda$.  But other regulators are possible, and some will in fact preserve
the Ward identity.  In our computations of the vertex and electron
self-energy diagrams, we used a Pauli--Villars regulator.  This regulator
preserved the relation $Z_1 = Z_2$, a consequence of the Ward identity; a
naive cutoff does not (see Problem~7.2).  We could fix our present
computation by introducing Pauli--Villars fermions.  Unfortunately, several
sets of fermions are required, making the method rather complicated.%
\footnote{This method is presented in Bjorken and Drell (1964), p.~154.}  We
will use a simpler method, \emph{dimensional regularization}, due to 't Hooft
and Veltman.%
\footnote{G.~'t~Hooft and M.~J.~G.~Veltman, \emph{Nucl.\ Phys.} \textbf{B44},
189 (1972).}  Dimensional regularization preserves the symmetries of QED and
also of a wide class of more general theories.

The question of which regulator to use has no a priori answer in quantum
field theory.  Often the choice has no effect on the predictions of the
theory.  When two regulators give different answers for observable
quantities, it is generally because some symmetry (such as the Ward identity)
is being violated by one (or both) of them.  In these cases we take the
symmetry to be fundamental and demand that it be preserved by the regulator.
But the validity of this choice cannot be proven; we are adopting the
symmetry as a new axiom.

\subsection{Dimensional Regularization}

The idea of dimensional regularization is very simple to state: Compute the
Feynman diagram as an analytic function of the dimensionality of spacetime,
$d$.  For sufficiently small $d$, any loop-momentum integral will converge and
therefore the Ward identity can be proved.  The final expression for any
observable quantity should have a well-defined limit as $d\to 4$.

Let us do a practice calculation to see how this technique works.  We consider
spacetime to have one time dimension and $(d-1)$ space dimensions.  Then we
can Wick-rotate Feynman integrals as before, to give integrals over a
$d$-dimensional Euclidean space.  A typical example is
\begin{equation}
  \int \frac{d^d\ell_E}{(2\pi)^d}\ \frac{1}{(\ell_E^2 + \Delta)^2}
  = \int_0^{\infty} \frac{d\ell_E}{(2\pi)^d}\ \frac{S_d\,\ell_E^{d-1}}
    {(\ell_E^2 + \Delta)^2} ,
  \label{eq:7.80}
\end{equation}
where $S_d = 2\pi^{d/2}/\Gamma(d/2)$ is the area of a unit sphere in $d$
dimensions.  To compute it, use the following trick:
\begin{equation}
  (\sqrt{\pi})^d = \Bigl( \int dx\ e^{-x^2} \Bigr)^d
    = \int d^d x\ e^{-|x|^2}
    = S_d \int_0^{\infty} dr\ r^{d-1} e^{-r^2}
    = S_d\cdot \tfrac{1}{2}\Gamma(d/2) .
\end{equation}
Evaluating the integral over $\ell_E$ using the standard formula
\begin{equation}
  \int_0^{\infty} d\ell\ \frac{\ell^{d-1}}{(\ell^2+\Delta)^2}
    = \Delta^{d/2-2}\ \frac{\Gamma(d/2)\Gamma(2-d/2)}{2},
\end{equation}
we obtain
\begin{equation}
  \int \frac{d^d\ell_E}{(2\pi)^d}\ \frac{1}{(\ell_E^2+\Delta)^2}
  = \frac{\Gamma(2-d/2)}{(4\pi)^{d/2}}\ \Delta^{d/2-2} .
  \label{eq:7.81}
\end{equation}
As $d\to 4$, the gamma function $\Gamma(2-d/2)$ has a pole.  Expanding about
$d=4$ with $d = 4-2\epsilon$, we use
\begin{equation}
  \Gamma(\epsilon) = \frac{1}{\epsilon} - \gamma_E + \mathcal{O}(\epsilon),
\end{equation}
where $\gamma_E = 0.5772\ldots$ is the Euler--Mascheroni constant.  The
divergence shows up as a pole in $\epsilon$, which will be absorbed by
renormalization.

Applying dimensional regularization to the vacuum polarization integral
\eqref{eq:7.79}, and contracting with $(q^2g^{\mu\nu} - q^\mu q^\nu)$, one
finds
\begin{equation}
  \Pi_2(q^2) = -\frac{2\alpha}{\pi}\ \int_0^1 dx\
    x(1-x)\ \Bigl[ \frac{1}{\epsilon} - \gamma_E + \log\Bigl(
      \frac{4\pi\mu^2}{\Delta} \Bigr) + \cdots \Bigr],
  \label{eq:7.82}
\end{equation}
where $\mu$ is the mass scale introduced by dimensional regularization.  The
physical, renormalized result is obtained by subtracting the value at
$q^2 = 0$:
\begin{equation}
  \Pi_2(q^2) - \Pi_2(0) = -\frac{2\alpha}{\pi}\ \int_0^1 dx\
    x(1-x)\ \log\Bigl( \frac{m^2}{m^2 - x(1-x)q^2} \Bigr) .
  \label{eq:7.91}
\end{equation}
Carrying out the integral,
\begin{equation}
  \Pi_2(q^2) - \Pi_2(0) = -\frac{2\alpha}{\pi}\ \Bigl[
    \frac{5}{9} - \frac{4m^2}{3q^2} + \frac{(4m^2/q^2)(1+2m^2/q^2)}
      {\sqrt{1-4m^2/q^2}}\ \mathrm{arctanh}\Bigl(
        \frac{1}{\sqrt{1-4m^2/q^2}} \Bigr) \Bigr],
  \label{eq:7.91b}
\end{equation}
for $q^2 < 0$, continued analytically to the physical region.

Let us calculate the imaginary part of $\Pi_2$ for $q^2 > 4m^2$.  For any
fixed $q^2$, the $x$-values that contribute are between the points
$x = \tfrac{1}{2}\pm\tfrac{1}{2}\beta$, where
$\beta = \sqrt{1-4m^2/q^2}$.  Since
$\mathrm{Im}\,[\log(-X\pm i\epsilon)] = \pm\pi$, we have
\begin{equation}
  \begin{aligned}
    \mathrm{Im}\,[\Pi_2(q^2\pm i\epsilon)]
    &= \mp \frac{2\alpha}{\pi}\cdot\pi\ \int_{\frac{1}{2}-\frac{\beta}{2}}
      ^{\frac{1}{2}+\frac{\beta}{2}} dx\ x(1-x)\\
    &= \mp 2\alpha\ \int_{-\beta/2}^{\beta/2} dy\ \Bigl(
      \frac{1}{4} - y^2 \Bigr)
    = \mp \frac{2\alpha}{3}\ \beta\Bigl( 1 - \frac{\beta^2}{4} \Bigr)\ \cdots
    = \mp \frac{\alpha}{3}\ \beta(1 - \tfrac{1}{3}\beta^2)\ \cdots .
  \end{aligned}
  \label{eq:7.92}
\end{equation}
This dependence on $q^2$ is exactly the same as in Eq.~\eqref{eq:5.13}, the
cross section for production of a fermion-antifermion pair.  That is just what
we would expect from the unitarity relation shown in Fig.~\ref{fig:7.6}(b);
the cut through the diagram for forward Bhabha scattering gives the total
cross section for $e^{+}e^{-}\to f\bar{f}$.  The parameter $\beta$ is
precisely the velocity of the fermions in the center-of-mass frame.

\begin{figure}[t]
  \centering
  \includegraphics[width=0.6\textwidth]{fig7.6.png}
  \caption{Two contributions to the optical theorem for Bhabha scattering.}
  \label{fig:7.6}
\end{figure}

Next let us examine how $\Pi_2(q^2)$ modifies the electromagnetic
interaction, as determined by Eq.~\eqref{eq:7.77}.  In the nonrelativistic
limit it makes sense to compute the potential $V(\mathbf{r})$.  For the
interaction between unlike charges, we have, in analogy with Eq.~\eqref{eq:4.135},
\begin{equation}
  V(\mathbf{x}) = \int \frac{d^3q}{(2\pi)^3}\ e^{i\mathbf{q}\cdot\mathbf{x}}\
    \frac{-e^2}{|\mathbf{q}|^2\,[1 - \Pi_2(-|\mathbf{q}|^2)]} .
  \label{eq:7.93}
\end{equation}
Expanding $\Pi_2$ for $|q^2| \ll m^2$, we obtain
\begin{equation}
  V(\mathbf{r}) = -\frac{\alpha}{r}\ \Bigl[ 1 + \frac{4\alpha}{5\pi}\,
    \frac{e^{-2mr}}{(mr)^{3/2}}\ \cdots \Bigr]
  \approx -\frac{\alpha}{r}\ -\ \frac{4\alpha^2}{5\sqrt{\pi}m^{3/2}}\,
    \delta^{(3)}(\mathbf{x}) .
  \label{eq:7.94}
\end{equation}
The correction term indicates that the electromagnetic force becomes much
stronger at small distances.  This effect can be measured in the hydrogen
atom, where the energy levels are shifted by
\begin{equation}
  \Delta E = \int d^3x\ |\psi(\mathbf{x})|^2\ \Bigl(
    -\frac{4\alpha^2}{5\sqrt{\pi}m^{3/2}}\delta^{(3)}(\mathbf{x}) \Bigr)
  = -\frac{4\alpha^2}{5\sqrt{\pi}m^{3/2}}\ |\psi(0)|^2 .
\end{equation}
The wavefunction $\psi(\mathbf{x})$ is nonzero at the origin only for $s$-wave
states.  For the $2S$ state, the shift is
\begin{equation}
  \Delta E = -\frac{4\alpha^2}{5\sqrt{\pi}m^{3/2}}\ \frac{1}{\pi a_0^3}\ \cdots
  = -\frac{\alpha^5 m}{15\sqrt{2}\pi}\ \cdots .
\end{equation}
This is a (small) part of the Lamb shift splitting listed in
Table~\ref{tab:6.1}.

\begin{figure}[t]
  \centering
  \includegraphics[width=0.5\textwidth]{fig7.7.png}
  \caption{Contour for evaluating the effective strength of the
    electromagnetic interaction in the nonrelativistic limit.  The pole at
    $Q = i\mu$ gives the Coulomb potential.  The branch cut gives the
    order-$\alpha$ correction due to vacuum polarization.}
  \label{fig:7.7}
\end{figure}

The delta function in Eq.~\eqref{eq:7.94} is only an approximation; to find
the true range of the correction term, we write Eq.~\eqref{eq:7.93} in the
form
\begin{equation}
  V(\mathbf{r}) = \int \frac{d^3q}{(2\pi)^3}\ e^{i\mathbf{q}\cdot\mathbf{r}}\
    \frac{-\alpha}{|\mathbf{q}|^2 + \mu^2}\ \Bigl[ 1 +
      \Pi_2(-|\mathbf{q}|^2) + \cdots \Bigr],
\end{equation}
where we have inserted a photon mass $\mu$ to regulate the Coulomb potential.
To perform this integral we push the contour upward (see
Fig.~\ref{fig:7.7}).  The leading contribution comes from the pole at
$Q = i\mu$, giving the Coulomb potential, $-\alpha/r$.  But there is an
additional contribution from the branch cut, which begins at
$Q = 2mi$.  The real part of the integrand is the same on both sides of the
cut, so the only contribution to the integral comes from the imaginary part
of $\Pi_2$.  Defining $q = -iQ$, we find that the contribution from the cut is
\begin{equation}
  \begin{aligned}
    \delta V(\mathbf{r}) &= \frac{-\alpha}{\pi^2 r}\ \int_{2m}^{\infty} dq\
      \frac{\mathrm{Im}\,\Pi_2(-q^2)}{q^2}\ \sin(qr)\ \cdots
    = \frac{-\alpha}{\pi r}\ \int_{2m}^{\infty} dq\
      \frac{2\alpha}{3}\ \frac{\beta(1-\tfrac{1}{3}\beta^2)}{q^2}\ \sin(qr)\ \cdots .
  \end{aligned}
\end{equation}
When $r \ll 1/m$, this integral is dominated by the region where
$q\sim 2m$.  Approximating the integrand in this region and substituting
$t = q - 2m$, we find
\begin{equation}
  \delta V(\mathbf{r}) = \frac{-\alpha}{r}\ \frac{4\alpha}{\sqrt{\pi}}\
    \int_0^{\infty} dt\ \frac{e^{-(t+2m)r}}{(mr)^{3/2}}\ \cdots
  = -\frac{\alpha}{r}\ \frac{4\alpha}{5\sqrt{\pi}}\,
    \frac{e^{-2mr}}{(mr)^{3/2}} ,
  \label{eq:7.95}
\end{equation}
so that
\begin{equation}
  V(\mathbf{r}) = -\frac{\alpha}{r}\ \Bigl[ 1 + \frac{4\alpha}{5\sqrt{\pi}}\,
    \frac{e^{-2mr}}{(mr)^{3/2}} \Bigr] .
  \label{eq:7.95b}
\end{equation}
Thus the range of the correction term is roughly the electron Compton
wavelength, $1/m$.  Since hydrogen wavefunctions are nearly constant on this
scale, the delta function in Eq.~\eqref{eq:7.94} was a good approximation.
The radiative correction to $V(\mathbf{r})$ is called the \emph{Uehling
potential}.

\begin{figure}[t]
  \centering
  \includegraphics[width=0.5\textwidth]{fig7.8.png}
  \caption{Virtual $e^{+}e^{-}$ pairs are effectively dipoles of length
    $\sim 1/m$, which screen the bare charge of the electron.}
  \label{fig:7.8}
\end{figure}

We can interpret the correction as being due to screening.  At $r\gg 1/m$,
virtual $e^{+}e^{-}$ pairs make the vacuum a dielectric medium in which the
apparent charge is less than the true charge (see Fig.~\ref{fig:7.8}).  At
smaller distances we begin to penetrate the polarization cloud and see the
bare charge.  This phenomenon is known as \emph{vacuum polarization}.

Now consider the opposite limit: small distance or $-q^2 \gg m^2$.
Equation \eqref{eq:7.91} then becomes
\begin{equation}
  \Pi_2(q^2) \approx \frac{2\alpha}{\pi}\ \int_0^1 dx\
    x(1-x)\ \Bigl[ \log\Bigl( \frac{-q^2}{m^2} \Bigr)
      + \log(x(1-x)) + \mathcal{O}\Bigl( \frac{m^2}{q^2} \Bigr) \Bigr] .
\end{equation}
The effective coupling constant in this limit is therefore
\begin{equation}
  \alpha_{\mathrm{eff}}(q^2) = \frac{\alpha}{1 +
    \frac{2\alpha}{3\pi}\log\Bigl( \frac{-q^2}{m^2} \Bigr)
      + \cdots }
  \xrightarrow{-q^2 \gg m^2}
    \frac{\alpha}{1 - \frac{\alpha}{3\pi}\log\Bigl(
      \frac{-q^2}{\Lambda^2} \Bigr)} ,
  \label{eq:7.96}
\end{equation}
where $\Lambda = \exp(5/3)\,m$.  The effective electric charge becomes much
larger at small distances, as we penetrate the screening cloud of virtual
electron-positron pairs.

\begin{figure}[t]
  \centering
  \includegraphics[width=0.7\textwidth]{fig7.9.png}
  \caption{Differential cross section for Bhabha scattering,
    $e^{+}e^{-}\to e^{+}e^{-}$, at $E_{\mathrm{cm}} = 29$ GeV, as measured by
    the HRS collaboration, M.~Derrick, et.\ al., \emph{Phys.\ Rev.\ D}
    \textbf{34}, 3286 (1986).  The upper curve is the order-$\alpha^2$
    prediction derived in Problem~5.2, plus a very small ($\sim 2\%$)
    correction due to the weak interaction.  The lower curve includes all QED
    radiative corrections to order $\alpha^3$ except the vacuum polarization
    contribution; note that these corrections depend on the experimental
    conditions, as explained in Chapter~6.  The middle curve includes the
    vacuum polarization contribution as well, which increases the effective
    value of $\alpha^2$ by about $10\%$ at this energy.}
  \label{fig:7.9}
\end{figure}

The combined vacuum polarization effect of the electron and of heavier quarks
and leptons causes the value of $\alpha_{\mathrm{eff}}(q^2)$ to increase by
about $5\%$ from $q = 0$ to $q = 30$ GeV, and this effect is observed in
high-energy experiments.  Figure~\ref{fig:7.9} shows the cross section for
Bhabha scattering at $E_{\mathrm{cm}} = 29$ GeV, and a comparison to QED with
and without the vacuum polarization diagram.  We can write
$\alpha_{\mathrm{eff}}$ as a function of $r$ by setting $q = 1/r$.  The
behavior of $\alpha_{\mathrm{eff}}(r)$ for all $r$ is sketched in
Fig.~\ref{fig:7.10}.  The idea of a distance-dependent (or
``scale-dependent'' or ``running'') coupling constant will be a major theme of
the rest of this book.

\begin{figure}[t]
  \centering
  \includegraphics[width=0.6\textwidth]{fig7.10.png}
  \caption{A qualitative sketch of the effective electromagnetic coupling
    constant generated by the one-loop vacuum polarization diagram, as a
    function of distance.  The horizontal scale covers many orders of
    magnitude.}
  \label{fig:7.10}
\end{figure}

\problems

\textbf{7.1} In Section~7.3 we used an indirect method to analyze the one-loop
$s$-channel diagram for boson-boson scattering in $\phi^4$ theory.  To verify
our indirect analysis, evaluate all three one-loop diagrams, using the
standard method of Feynman parameters.  Check the validity of the optical
theorem.

\textbf{7.2} \emph{Alternative regulators in QED.}  In Section~7.5, we saw
that the Ward identity can be violated by an improperly chosen regulator.
Let us check the validity of the identity $Z_1 = Z_2$, to order $\alpha$, for
several choices of the regulator.  We have already verified that the relation
holds for Pauli--Villars regularization.
\begin{enumerate}[label=(\alph*)]
  \item Recompute $\delta Z_1$ and $\delta Z_2$, defining the integrals
    Eq.~(6.49) and Eq.~(6.50) by simply placing an upper limit
    $\Lambda$ on the integration over $|\ell_E|$.  Show that, with this
    definition, $\delta Z_1 \neq \delta Z_2$.
  \item Recompute $\delta Z_1$ and $\delta Z_2$, defining the integrals
    Eq.~(6.49) and Eq.~(6.50) by dimensional regularization.  You
    may take the Dirac matrices to be $4\times4$ as usual, but note that, in
    $d$ dimensions,
    \begin{equation}
      g^\mu_{\ \mu}\gamma^\nu\gamma_\nu = d .
    \end{equation}
    Show that, with this definition, $\delta Z_1 = \delta Z_2$.
\end{enumerate}

\textbf{7.3} In the text we discussed the Ward identity for amplitudes
involving an external photon.  There is an important generalization to
amplitudes involving \emph{two} external photons, known as the
Ward--Takahashi identity for the photon-photon amplitude.  Consider the
amplitude $\mathcal{M}^{\mu\nu}(k_1, k_2)$ for the production of two photons
with momenta $k_1$ and $k_2$.  Show that the amplitude satisfies
\begin{equation}
  k_{1\mu}\mathcal{M}^{\mu\nu}(k_1, k_2) = 0 .
\end{equation}
(You may want to use the result of the text that the sum over insertions of a
photon along a through-going fermion line telescopes, and that the sum over
insertions along a closed loop vanishes.)

\textbf{7.4} Verify that the result of the text for the vacuum polarization,
Eq.~\eqref{eq:7.91}, satisfies the unitarity relation
\begin{equation}
  \mathrm{Im}\,\Pi_2(q^2) = \frac{1}{2}\ \int d\Pi_2\
    \mathcal{M}(e^{+}e^{-}\to f\bar{f})\ ,
\end{equation}
in the physical region $q^2 > 4m^2$.

\textbf{7.5} In this problem we will derive the relation between the field
strength renormalization constant $Z$ and the physical propagator pole that
was quoted in Section~7.1.
\begin{enumerate}[label=(\alph*)]
  \item Consider the exact two-point function $D(p^2)$, and write down its
    K\"all\'en--Lehmann spectral representation.
  \item Show that the residue of the pole at $p^2 = m^2$ is $iZ$, by
    expanding the spectral representation near the pole.
  \item Show that the renormalized propagator, defined by $D_R(p^2) =
    Z^{-1}D(p^2)$, has unit residue at the pole.
\end{enumerate}
